\documentclass[11pt,a4paper]{article}

% ------------------------------------------------------------
%  SOLUCIONARI DEL DOCENT
%  Per activar-lo, canvia \solucionsfalse per \solucionstrue
% ------------------------------------------------------------
\newif\ifsolucions
\solucionsfalse
\ifdefined\solucionsactives\solucionstrue\fi

% ============================================================
%  PREÀMBUL — Dossier 1: Circuits amb resistències (4t ESO)
% ============================================================
\usepackage[utf8]{inputenc}
\usepackage[T1]{fontenc}
\usepackage{textcomp}
% El punt volat (·) es dibuixa amb el símbol matemàtic \cdot: així es veu igual
% a qualsevol tipus de lletra i amb qualsevol motor (pdfLaTeX, XeTeX/Tectonic, LuaTeX).
\usepackage{newunicodechar}
\newunicodechar{·}{\ensuremath{\cdot}}
\newcommand{\sep}{\ \ensuremath{\cdot}\ }
\usepackage[catalan]{babel}
\usepackage[sc]{mathpazo}
\usepackage[a4paper,left=2cm,right=2cm,top=2.3cm,bottom=2.2cm]{geometry}
\usepackage{microtype}
\usepackage{xcolor}
\usepackage{amsmath,amssymb}
\usepackage{siunitx}
\sisetup{output-decimal-marker={,},group-minimum-digits=5,group-separator={\,},
         detect-all=false,per-mode=symbol}
% Amb XeTeX/LuaTeX (p. ex. Tectonic), el símbol Ω que dona siunitx no existeix a la lletra:
% el dibuixem amb la lletra grega de les matemàtiques.
\usepackage{iftex}
\ifPDFTeX\else\DeclareSIUnit{\ohm}{\ensuremath{\mathrm{\Omega}}}\fi
\usepackage[european]{circuitikz}
\usepackage[most]{tcolorbox}
\usepackage{booktabs,array,enumitem,fancyhdr,titlesec,needspace,pifont,graphicx}
\usepackage[hidelinks]{hyperref}

\linespread{1.08}
\setlength{\parindent}{0pt}
\setlength{\parskip}{0.5em}
\setlist{nosep,leftmargin=1.4em}

% ---------- Colors ----------
\definecolor{blau}{HTML}{1F5F99}
\definecolor{verd}{HTML}{2E7D32}
\definecolor{taronja}{HTML}{E07B00}
\definecolor{vermell}{HTML}{B3261E}
\definecolor{lila}{HTML}{6A3FA0}

% ---------- Títols ----------
\titleformat{\section}{\Large\bfseries\color{blau}}{Bloc \thesection}{0.7em}{}[\vspace{-0.4em}\color{blau}\titlerule]
\titleformat{\subsection}{\large\bfseries\color{blau}}{}{0em}{}
\titlespacing*{\section}{0pt}{0pt}{1em}
\titlespacing*{\subsection}{0pt}{1.2em}{0.4em}

% ---------- Capçalera i peu ----------
\pagestyle{fancy}
\fancyhf{}
\renewcommand{\headrulewidth}{0.3pt}
\fancyhead[L]{\small\textsc{Electrònica analògica}\sep 4t ESO}
\fancyhead[R]{\small Dossier 1\sep Circuits amb resistències\ifsolucions\sep\textbf{SOLUCIONARI}\fi}
\fancyfoot[C]{\small\thepage}

% ---------- Unitats i símbols de circuit ----------
\ctikzset{bipoles/length=1.25cm, bipoles/resistor/height=0.28, bipoles/resistor/width=0.7}

% ---------- Quadres de colors ----------
\tcbset{boxrule=0.6pt, arc=3pt, left=7pt, right=7pt, top=5pt, bottom=5pt,
        fonttitle=\bfseries, breakable, before skip=0.8em, after skip=0.8em}
\newtcolorbox{idea}{colback=blau!6, colframe=blau, title={La idea clau}}
\newtcolorbox{metode}[1]{colback=lila!6, colframe=lila, title={Mètode: #1}}
\newtcolorbox{exemple}[1]{colback=verd!7, colframe=verd, title={Exemple resolt: #1}}
\newtcolorbox{guiada}[1]{colback=taronja!8, colframe=taronja, title={Practiquem junts: #1}}
\newtcolorbox{atencio}{colback=vermell!6, colframe=vermell, title={Compte amb l'error típic}}
\newtcolorbox{comprova}{colback=yellow!12, colframe=yellow!55!black, title={Comprova-ho}}

% ---------- Nivells d'exercici ----------
\newcommand{\estrella}{\ding{72}}
\newcommand{\nivell}[1]{%
  \ifcase#1\or{\color{verd}\estrella}\or{\color{taronja}\estrella\estrella}\or{\color{vermell}\estrella\estrella\estrella}\fi}

\newcounter{exercici}[section]
\renewcommand{\theexercici}{\thesection.\arabic{exercici}}
\newenvironment{exercici}[2][8]{%
  \par\bigskip\needspace{#1\baselineskip}\noindent
  \refstepcounter{exercici}\textbf{\theexercici}\enspace\nivell{#2}\enspace\ignorespaces}%
  {\par}

% ---------- Solucions (espai en blanc o quadre de solució) ----------
%  \resposta{espai en blanc}{solució}
\newcommand{\resposta}[2]{%
  \ifsolucions
    \par\smallskip
    \begin{tcolorbox}[colback=verd!6, colframe=verd!60!black, boxrule=0.5pt, arc=2pt,
                      left=5pt, right=5pt, top=3pt, bottom=3pt, before skip=0.3em, after skip=0.6em]
      {\small\textbf{Solució.}\ #2}
    \end{tcolorbox}
  \else
    \vspace*{#1}
  \fi}

% ---------- Passos del mètode (per a les pràctiques guiades) ----------
\newcommand{\pas}[3]{\par\noindent\textbf{\textcolor{taronja}{\ding{\numexpr171+#1\relax}~#2}}\quad #3\par\smallskip}

% ---------- Taula de magnituds ----------
\newenvironment{taulamag}{%
  \par\smallskip\renewcommand{\arraystretch}{1.7}\centering
  \begin{tabular}{@{}l *{4}{>{\centering\arraybackslash}p{2.5cm}}@{}}
  \toprule
  \textbf{Element} & $V$ (\si{\volt}) & $I$ (\si{\ampere}) & $R$ (\si{\ohm}) & $P$ (\si{\watt})\\
  \midrule}%
  {\bottomrule\end{tabular}\par\smallskip}

% ---------- Dibuixos de circuits ----------
\newcommand{\simbol}[1]{\begin{circuitikz}[scale=0.8,transform shape]\draw (0,0) to[#1] (1.6,0);\end{circuitikz}}

% Una resistència: V, R
\newcommand{\csimple}[2]{%
\begin{circuitikz}[scale=0.9,transform shape]
 \draw (0,0) to[battery1,l_={#1}] (0,2.2) to[R,l={#2}] (3.5,2.2) to[short] (3.5,0) to[short] (0,0);
\end{circuitikz}}

% Dues resistències en sèrie
\newcommand{\cserie}[3]{%
\begin{circuitikz}[scale=0.9,transform shape]
 \draw (0,0) to[battery1,l_={#1}] (0,2.2) to[R,l={#2}] (3.2,2.2) to[R,l={#3}] (6.4,2.2)
       to[short] (6.4,0) to[short] (0,0);
\end{circuitikz}}

% Tres resistències en sèrie
\newcommand{\cserietres}[4]{%
\begin{circuitikz}[scale=0.8,transform shape]
 \draw (0,0) to[battery1,l_={#1}] (0,2.2) to[R,l={#2}] (3,2.2) to[R,l={#3}] (6,2.2) to[R,l={#4}] (9,2.2)
       to[short] (9,0) to[short] (0,0);
\end{circuitikz}}

% Dues resistències en paral·lel
\newcommand{\cparalel}[3]{%
\begin{circuitikz}[scale=0.9,transform shape]
 \draw (0,0) to[battery1,l_={#1}] (0,2.6) to[short] (3,2.6) to[R,l={#2}] (3,0) to[short] (0,0);
 \draw (3,2.6) to[short,*-] (6,2.6) to[R,l={#3}] (6,0) to[short,-*] (3,0);
\end{circuitikz}}

% Tres resistències en paral·lel
\newcommand{\cparaleltres}[4]{%
\begin{circuitikz}[scale=0.85,transform shape]
 \draw (0,0) to[battery1,l_={#1}] (0,2.6) to[short] (3,2.6) to[R,l={#2}] (3,0) to[short] (0,0);
 \draw (3,2.6) to[short,*-] (6,2.6) to[R,l={#3}] (6,0) to[short,-*] (3,0);
 \draw (6,2.6) to[short,*-] (9,2.6) to[R,l={#4}] (9,0) to[short,-*] (6,0);
\end{circuitikz}}

% Mixt A: R1 en sèrie amb (R2 en paral·lel amb R3)
\newcommand{\cmixa}[4]{%
\begin{circuitikz}[scale=0.85,transform shape]
 \draw (0,0) to[battery1,l_={#1}] (0,2.6) to[R,l={#2}] (3,2.6) to[short,*-] (6,2.6) to[R,l={#4}] (6,0)
       to[short,-*] (3,0) to[short] (0,0);
 \draw (3,2.6) to[R,l={#3}] (3,0);
\end{circuitikz}}

% Mixt C: (R1 en sèrie amb R2) en paral·lel amb R3
\newcommand{\cmixc}[4]{%
\begin{circuitikz}[scale=0.85,transform shape]
 \draw (0,0) to[battery1,l_={#1}] (0,2.6) to[short] (2,2.6) to[R,l={#4}] (2,0) to[short] (0,0);
 \draw (2,2.6) to[short,*-] (4,2.6) to[R,l={#2}] (7,2.6) to[R,l={#3}] (10,2.6) to[short] (10,0) to[short] (2,0);
 \node[circ] at (2,0) {};
\end{circuitikz}}

% Mixt D: R1 en sèrie amb (R2 en paral·lel amb R3) i amb R4
\newcommand{\cmixd}[5]{%
\begin{circuitikz}[scale=0.8,transform shape]
 \draw (0,0) to[battery1,l_={#1}] (0,2.6) to[R,l={#2}] (3,2.6) to[short,*-] (6,2.6) to[R,l={#4}] (6,0)
       to[short,-*] (3,0) to[R,l_={#5}] (0,0);
 \draw (3,2.6) to[R,l={#3}] (3,0);
\end{circuitikz}}


\begin{document}

% ---------------- Portada ----------------
\begin{titlepage}
\centering
\vspace*{2.5cm}
{\small\textsc{Tecnologia\sep 4t ESO}}\par\vspace{0.4cm}
{\Huge\bfseries\color{blau} Electrònica analògica}\par\vspace{0.3cm}
{\LARGE Dossier 1}\par\vspace{0.6cm}
{\Huge\bfseries Circuits amb resistències}\par\vspace{0.4cm}
{\Large\itshape pas a pas}\par\vspace{1.6cm}
\begin{circuitikz}[scale=1.1]
 \draw (0,0) to[battery1] (0,2) to[R] (3,2) to[R] (6,2) to[short] (6,0) to[short] (0,0);
\end{circuitikz}
\vfill
\begin{flushleft}\large
\textbf{Nom i cognoms:}\par\vspace{1.0cm}
\textbf{Grup:}\par\vspace{1.0cm}
\textbf{Data d'inici:}
\end{flushleft}
\ifsolucions\vspace{0.5cm}{\Large\bfseries\color{vermell} SOLUCIONARI DEL DOCENT}\fi
\end{titlepage}

% ---------------- Com funciona ----------------
\section*{Com fer servir aquest dossier}

Aquest dossier t'ajuda a aprendre a calcular circuits amb resistències \textbf{pas a pas}. No cal que t'ho sàpigues tot per començar: cada bloc parteix del que has vist abans.

\subsection*{Cada bloc segueix el mateix ritme}
\begin{enumerate}[label=\textbf{\arabic*.}]
 \item \textbf{Llegeix la idea clau.} És el més important, en poques línies.
 \item \textbf{Mira l'exemple resolt.} Fixa't en \emph{com} es fa cada pas.
 \item \textbf{Practica amb nosaltres.} Una part ja està feta i tu acabes la resta.
 \item \textbf{Fes els exercicis.} Comença pel nivell que et vagi bé.
 \item \textbf{Comprova-ho.} Una pregunta curta per veure si ho has entès.
\end{enumerate}

\subsection*{Els tres nivells}
\begin{center}
\renewcommand{\arraystretch}{1.5}
\begin{tabular}{@{}cll@{}}
\nivell{1} & \textbf{Nivell 1} & Ho he vist fer: aplico el mateix mètode amb altres números.\\
\nivell{2} & \textbf{Nivell 2} & Ho sé aplicar: hi ha més dades o més passos.\\
\nivell{3} & \textbf{Nivell 3} & Repte: haig de pensar quin camí fer servir.
\end{tabular}
\end{center}

\subsection*{Tres consells}
\begin{itemize}
 \item \textbf{Passa-ho tot a unitats del SI} (V, A, $\Omega$) abans de calcular.
 \item \textbf{Escriu sempre la fórmula} abans de posar-hi els números.
 \item \textbf{Comprova el resultat}: té sentit? Les tensions i els corrents sumen bé?
\end{itemize}

\clearpage
\tableofcontents
\clearpage

% ---------------- Blocs ----------------
\setcounter{section}{-1}
\section{Preparem-nos}

Abans de calcular circuits necessitem tres coses: saber \textbf{què és un circuit}, conèixer les \textbf{magnituds i les unitats} i dominar els \textbf{prefixos} (mil·li, quilo\ldots). Si això va bé, la resta del dossier és molt més fàcil.

\subsection{El circuit elèctric}

\begin{idea}
Un \textbf{circuit elèctric} és un camí \textbf{tancat} pel qual circula el corrent. Té, com a mínim:
\begin{itemize}
 \item un \textbf{generador} (una pila o una font), que dona l'energia;
 \item un \textbf{receptor} (una bombeta, un motor, una resistència\ldots), que l'aprofita;
 \item uns \textbf{conductors} (els cables), que ho connecten tot.
\end{itemize}
Si el circuit s'obre (per exemple, amb un interruptor), el corrent s'atura.
\end{idea}

\begin{center}\csimple{\SI{12}{V}}{$R=\SI{60}{\ohm}$}\end{center}

Els circuits es dibuixen amb \textbf{símbols}. Fem servir els símbols europeus, en què la resistència és un rectangle:

\begin{center}
\renewcommand{\arraystretch}{1.1}
\begin{tabular}{@{}c l@{\hspace{1.6cm}}c l@{}}
\simbol{battery1} & Generador (pila o font) & \simbol{R} & Resistència\\
\simbol{nos} & Interruptor & \simbol{lamp} & Bombeta\\
\simbol{ammeter} & Amperímetre & \simbol{voltmeter} & Voltímetre
\end{tabular}
\end{center}

\subsection{Magnituds i unitats}

\renewcommand{\arraystretch}{1.5}
\begin{center}
\begin{tabular}{@{}llllp{5.6cm}@{}}
\toprule
\textbf{Magnitud} & \textbf{Símbol} & \textbf{Unitat} & \textbf{Aparell} & \textbf{Pensa-hi com\ldots}\\
\midrule
Tensió (voltatge) & $V$ & volt (\si{\volt}) & voltímetre & l'\emph{empenta} de la pila, com la pressió d'una bomba d'aigua\\
Intensitat & $I$ & ampere (\si{\ampere}) & amperímetre & el \emph{cabal}: quanta càrrega passa cada segon\\
Resistència & $R$ & ohm (\si{\ohm}) & òhmmetre & l'\emph{estretor} d'un tub: dificulta el pas\\
Potència & $P$ & watt (\si{\watt}) & --- & l'energia que gasta el receptor \emph{cada segon}\\
\bottomrule
\end{tabular}
\end{center}

\subsection{Prefixos i conversions}

Les magnituds sovint surten amb prefixos: una pila pot donar \SI{250}{\milli\ampere}, i una resistència pot ser de \SI{4,7}{\kilo\ohm}. Cal saber passar-ho a \textbf{unitats base} abans de calcular.

\begin{center}
\renewcommand{\arraystretch}{1.4}
\begin{tabular}{@{}llll@{}}
\toprule
\textbf{Prefix} & \textbf{Símbol} & \textbf{Vol dir} & \textbf{Exemple}\\
\midrule
micro & $\mu$ & milionèsima ($\times 10^{-6}$) & \SI{8000}{\micro\ampere}\\
mil·li & m & mil·lèsima ($\times 10^{-3}$) & \SI{250}{\milli\ampere}\\
(unitat base) & --- & $\times 1$ & \SI{3}{\ampere}\\
quilo & k & mil vegades ($\times 10^3$) & \SI{4,7}{\kilo\ohm}\\
mega & M & un milió de vegades ($\times  10^6 $) & \SI{2,2}{\mega\ohm}\\
\bottomrule
\end{tabular}
\end{center}

\begin{metode}{passar a unitat base}
\begin{itemize}
 \item \textbf{\ding{172}} Mira el prefix. 
 \item \textbf{\ding{173}} Si és \emph{petit} (mil·li, micro), \textbf{divideix}: $\div 1000$ per mil·li, $\div 1\,000\,000$ per micro.\quad
 \item \textbf{\ding{174}} Si és \emph{gran} (quilo, mega), \textbf{multiplica}: $\times 1000$ per quilo, $\times 1\,000\,000$ per mega.
\end{itemize}
\end{metode}

\begin{exemple}{conversions}
\begin{itemize}
 \item \SI{470}{\milli\ampere} $= 470 \times 10^{-3} =$ \SI{4,7e-1}{\ampere}
 \item \SI{2,2}{\kilo\ohm} $= 2{,}2 \times 10^3 =$ \SI{2,2e3}{\ohm}
 \item \SI{0,05}{\ampere} $= 5 \times 10^{-3} \text{ A} =$ \SI{50}{\milli\ampere}\quad (a l'inrevés: de base a mil·li, \emph{multipliquem})
 \item \SI{1}{\mega\ohm} $= 1 \times 10^6 \ \Omega =$ \SI{1e6}{\ohm}
\end{itemize}
\end{exemple}

\begin{atencio}
No confonguis \textbf{m} (mil·li, una mil·lèsima) amb \textbf{M} (mega, un milió). \SI{1}{\milli\ampere} és molt petit; \SI{1}{\mega\ohm} és molt gran.
\end{atencio}

\clearpage
\subsection*{Exercicis}

\begin{exercici}{1}
Utilitza els factors de conversió per passar a unitats base (A, V o $\Omega$):\quad
\textbf{a)} \SI{250}{\milli\ampere}\qquad \textbf{b)} \SI{3,3}{\kilo\ohm}\qquad \textbf{c)} \SI{12}{\milli\volt}
\resposta{2.2cm}{a) \SI{0,25}{\ampere}\quad b) \SI{3300}{\ohm}\quad c) \SI{0,012}{\volt}}
\end{exercici}

\begin{exercici}{1}
Passa a la unitat que s'indica:\quad
\textbf{a)} \SI{0,5}{\ampere} a mA\qquad \textbf{b)} \SI{4700}{\ohm} a k$\Omega$\qquad \textbf{c)} \SI{0,02}{\ampere} a mA
\resposta{2.2cm}{a) \SI{500}{\milli\ampere}\quad b) \SI{4,7}{\kilo\ohm}\quad c) \SI{20}{\milli\ampere}}
\end{exercici}

\begin{exercici}{2}
Fes aquestes conversions:\quad
\textbf{a)} \SI{2,2}{\mega\ohm} a k$\Omega$\qquad \textbf{b)} \SI{1500}{\milli\ampere} a A\qquad \textbf{c)} \SI{0,0047}{\kilo\ohm} a $\Omega$\qquad \textbf{d)} \SI{8000}{\micro\ampere} a mA
\resposta{2.2cm}{a) \SI{2200}{\kilo\ohm}\quad b) \SI{1,5}{\ampere}\quad c) \SI{4,7}{\ohm}\quad d) \SI{8}{\milli\ampere}}
\end{exercici}

\begin{exercici}{2}
Observa el circuit i respon.
\begin{center}
\begin{circuitikz}[scale=0.9,transform shape]
 \draw (0,0) to[battery1,l_=\SI{9}{V}] (0,2.2) to[nos,l=interruptor] (3.2,2.2) to[lamp,l=bombeta] (6.4,2.2)
       to[short] (6.4,0) to[short] (0,0);
\end{circuitikz}
\end{center}
\textbf{a)} Quin component fa de generador i quin de receptor?\quad \textbf{b)} Circula corrent ara mateix? Per què?
\resposta{2.2cm}{a) El generador és la pila i el receptor és la bombeta. b) No: l'interruptor està obert, per tant el circuit no és tancat i no hi circula corrent.}
\end{exercici}

\begin{exercici}{3}
Ordena de menor a major aquestes intensitats:\quad \SI{0,02}{\ampere},\ \ \SI{15}{\milli\ampere},\ \ \SI{8000}{\micro\ampere}.
\resposta{2.2cm}{Tot en mA: \SI{20}{\milli\ampere}, \SI{15}{\milli\ampere}, \SI{8}{\milli\ampere}. Ordre: \SI{8000}{\micro\ampere} $<$ \SI{15}{\milli\ampere} $<$ \SI{0,02}{\ampere}.}
\end{exercici}

\clearpage
\section{La llei d'Ohm}

Amb una pila, una resistència i un amperímetre, el físic Georg Ohm va descobrir una relació molt senzilla entre la tensió, la intensitat i la resistència. Aquesta relació és la base de tot el que farem després.

\begin{idea}
La \textbf{llei d'Ohm} diu que la tensió que hi ha en una resistència és igual al producte de la intensitat per la resistència:
\[ V = I \cdot R \]
\begin{itemize}
 \item Si la \textbf{tensió augmenta}, la intensitat també augmenta.
 \item Si la \textbf{resistència augmenta}, la intensitat \textbf{disminueix}.
\end{itemize}
\end{idea}

\begin{center}
\begin{minipage}{0.4\linewidth}\centering
\begin{tikzpicture}[scale=0.85]
 \draw[blau,thick] (0,0) -- (4,0) -- (2,3) -- cycle;
 \draw[blau,thick] (0.9,1.35) -- (3.1,1.35);
 \draw[blau,thick] (2,1.35) -- (2,0);
 \node at (2,2.0) {\Large $V$};
 \node at (1,0.6) {\Large $I$};
 \node at (3,0.6) {\Large $R$};
\end{tikzpicture}
\end{minipage}
\begin{minipage}{0.5\linewidth}
\textbf{El triangle de la llei d'Ohm.} Tapa la magnitud que busques amb el dit:
\begin{itemize}
 \item Tapes $V$: queda $I \cdot R$, és a dir $V = I \cdot R$.
 \item Tapes $I$: queda $V / R$, és a dir $I = \dfrac{V}{R}$.
 \item Tapes $R$: queda $V / I$, és a dir $R = \dfrac{V}{I}$.
\end{itemize}
\end{minipage}
\end{center}

\begin{metode}{resoldre un problema amb la llei d'Ohm}
\begin{itemize}
\item \textbf{\ding{172}} Escriu les \textbf{dades} amb unitats base (V, A, $\Omega$). \quad
\item \textbf{\ding{173}} Escriu la \textbf{fórmula} que necessites (amb el triangle). \quad
\item \textbf{\ding{174}} \textbf{Substitueix} i calcula. \quad
\item \textbf{\ding{175}} Escriu el \textbf{resultat amb la unitat} i pensa si té sentit.
\end{itemize}
\end{metode}

\begin{exemple}{calcular la intensitat}
Una resistència de \SI{60}{\ohm} es connecta a una pila de \SI{12}{V}. Quina intensitat hi passa?
\begin{itemize}
 \item \textbf{Dades:} $V = \SI{12}{V}$, \ $R = \SI{60}{\ohm}$. \ Busquem $I$.
 \item \textbf{Fórmula:} $I = \dfrac{V}{R}$
 \item \textbf{Substitució:} $I = \dfrac{12}{60} = \SI{0,2}{A} = \SI{200}{mA}$
\end{itemize}
\end{exemple}

\begin{exemple}{calcular la resistència}
Una bombeta connectada a \SI{9}{V} és travessada per \SI{30}{mA}. Quina resistència té?
\begin{itemize}
 \item \textbf{Dades:} $V = \SI{9}{V}$, \ $I = \SI{30}{mA} = \SI{0,03}{A}$. \ Busquem $R$.
 \item \textbf{Fórmula:} $R = \dfrac{V}{I}$
 \item \textbf{Substitució:} $R = \dfrac{9}{0{,}03} = \SI{300}{\ohm}$
\end{itemize}
\end{exemple}

\begin{exemple}{calcular la tensió}
Per una resistència de \SI{4,7}{\kilo\ohm} hi passa un corrent de \SI{2}{mA}. Quina tensió hi ha entre els seus extrems?
\begin{itemize}
 \item \textbf{Dades:} $R = \SI{4,7}{\kilo\ohm} = \SI{4700}{\ohm}$, \ $I = \SI{2}{mA} = \SI{0,002}{A}$. \ Busquem $V$.
 \item \textbf{Fórmula:} $V = I \cdot R$
 \item \textbf{Substitució:} $V = 0{,}002 \cdot 4700 = \SI{9,4}{V}$
\end{itemize}
\end{exemple}

\begin{atencio}
Oblidar passar els mA o els k$\Omega$ a unitats base. Si a l'exemple 2 fessis $9 / 30$, obtindries $0{,}3$: un resultat \textbf{1000 vegades} massa petit. Passa-ho \textbf{sempre} a V, A i $\Omega$ abans de calcular.
\end{atencio}

\begin{guiada}{una resistència de \SI{220}{\ohm} connectada a \SI{5}{V}}
Quina intensitat hi passa? Dóna el resultat en mA.
\pas{1}{Dades:}{$V = \SI{5}{V}$, \ $R = \SI{220}{\ohm}$. \ Busquem $I$.}
\pas{2}{Fórmula:}{\vspace{0.8cm}}
\pas{3}{Substitució:}{\vspace{0.8cm}}
\pas{4}{Resultat en mA:}{\vspace{0.8cm}}
\resposta{0pt}{$I = V/R = 5/220 \approx \SI{0,0227}{A} \approx \SI{22,7}{mA}$.}
\end{guiada}

\subsection*{Exercicis}

\begin{exercici}{1}
Calcula la intensitat d'una resistència de \SI{50}{\ohm} connectada a \SI{10}{V}.
\resposta{2.2cm}{$I = V/R = 10/50 = \SI{0,2}{A}$ (\SI{200}{mA}).}
\end{exercici}

\begin{exercici}{1}
Calcula la intensitat d'una resistència de \SI{120}{\ohm} connectada a \SI{6}{V}. Dóna el resultat en mA.
\resposta{2.2cm}{$I = 6/120 = \SI{0,05}{A} = \SI{50}{mA}$.}
\end{exercici}

\begin{exercici}{1}
Quina tensió hi ha entre els extrems d'una resistència de \SI{24}{\ohm} si hi passa \SI{0,5}{A}?
\resposta{2.2cm}{$V = I \cdot R = 0{,}5 \cdot 24 = \SI{12}{V}$.}
\end{exercici}

\begin{exercici}{2}
Una resistència connectada a \SI{9}{V} és travessada per \SI{45}{mA}. Quina resistència té?
\resposta{2.4cm}{$I = \SI{0,045}{A}$. $R = V/I = 9/0{,}045 = \SI{200}{\ohm}$.}
\end{exercici}

\begin{exercici}{2}
Per una resistència de \SI{2,2}{\kilo\ohm} hi passa un corrent de \SI{5}{mA}. Quina tensió hi ha entre els extrems?
\resposta{2.4cm}{$R = \SI{2200}{\ohm}$, $I = \SI{0,005}{A}$. $V = 0{,}005 \cdot 2200 = \SI{11}{V}$.}
\end{exercici}

\begin{exercici}{2}
Connectem una resistència de \SI{330}{\ohm} a una font de \SI{5}{V}. Quina intensitat hi passa? Dóna el resultat en mA.
\resposta{2.4cm}{$I = 5/330 \approx \SI{0,01515}{A} \approx \SI{15,2}{mA}$.}
\end{exercici}

\begin{exercici}{3}
Completa la taula aplicant la llei d'Ohm a cada fila:
\begin{center}
\renewcommand{\arraystretch}{1.9}
\begin{tabular}{@{}c *{3}{>{\centering\arraybackslash}p{2.8cm}}@{}}
\toprule
 & $V$ & $I$ & $R$\\
\midrule
a) & \SI{12}{V} & \SI{0,02}{A} & \\
b) & & \SI{3}{mA} & \SI{1}{\kilo\ohm}\\
c) & \SI{4,5}{V} & & \SI{1,5}{\kilo\ohm}\\
\bottomrule
\end{tabular}
\end{center}
\resposta{1cm}{a) $R = 12/0{,}02 = \SI{600}{\ohm}$. \quad b) $V = 0{,}003 \cdot 1000 = \SI{3}{V}$. \quad c) $I = 4{,}5/1500 = \SI{0,003}{A} = \SI{3}{mA}$.}
\end{exercici}

\begin{exercici}{3}
Una resistència connectada a \SI{12}{V} és travessada per \SI{0,3}{A}.
\textbf{a)} Quina resistència té?\quad \textbf{b)} Quina intensitat hi passarà si la connectem a \SI{6}{V}?\quad \textbf{c)} I si la connectem a \SI{36}{V}?\quad \textbf{d)} Què passa amb la intensitat quan la tensió es multiplica per un número?
\resposta{3cm}{a) $R = 12/0{,}3 = \SI{40}{\ohm}$. b) $I = 6/40 = \SI{0,15}{A}$ (la meitat). c) $I = 36/40 = \SI{0,9}{A}$ (el triple). d) La intensitat es multiplica pel mateix número: és \emph{proporcional} a la tensió.}
\end{exercici}

\begin{comprova}
Tenim una pila fixa i dupliquem la resistència del circuit. La intensitat es fa el doble o la meitat? Per què?
\resposta{2cm}{La meitat. Com que $I = V/R$, si $R$ es dobla i $V$ no canvia, la intensitat es divideix per dos.}
\end{comprova}

\clearpage
\section{La potència elèctrica}

Una bombeta de \SI{60}{W} il·lumina més que una de \SI{5}{W}. La \textbf{potència} ens diu quanta energia gasta un receptor \textbf{cada segon}. Saber-ho és important: una resistència amb massa potència s'escalfa i es pot cremar.

\begin{idea}
La \textbf{potència} d'un receptor és el producte de la tensió per la intensitat:
\[ P = V \cdot I \qquad \text{(en watts, W)} \]
Combinant-la amb la llei d'Ohm ($V = I \cdot R$) obtenim dues fórmules més:
\[ P = I^2 \cdot R \qquad\qquad P = \frac{V^2}{R} \]
Fes servir la que tingui les \textbf{dades que ja coneixes}.
\end{idea}

\begin{center}
\renewcommand{\arraystretch}{1.5}
\begin{tabular}{@{}lc@{}}
\toprule
\textbf{Si conec\ldots} & \textbf{Fórmula}\\
\midrule
$V$ i $I$ & $P = V \cdot I$\\
$I$ i $R$ & $P = I^2 \cdot R$\\
$V$ i $R$ & $P = V^2 / R$\\
\bottomrule
\end{tabular}
\end{center}

\begin{exemple}{una bombeta}
Una bombeta de \SI{12}{V} és travessada per \SI{0,5}{A}. Quina potència té?
\begin{itemize}
 \item \textbf{Dades:} $V = \SI{12}{V}$, \ $I = \SI{0,5}{A}$.
 \item \textbf{Fórmula:} $P = V \cdot I$
 \item \textbf{Resultat:} $P = 12 \cdot 0{,}5 = \SI{6}{W}$
\end{itemize}
\end{exemple}

\begin{exemple}{una resistència amb intensitat coneguda}
Per una resistència de \SI{100}{\ohm} hi passa un corrent de \SI{0,1}{A}. Quina potència dissipa?
\begin{itemize}
 \item \textbf{Dades:} $R = \SI{100}{\ohm}$, \ $I = \SI{0,1}{A}$.
 \item \textbf{Fórmula:} $P = I^2 \cdot R$
 \item \textbf{Resultat:} $P = 0{,}1^2 \cdot 100 = 0{,}01 \cdot 100 = \SI{1}{W}$
\end{itemize}
\end{exemple}

\begin{exemple}{aguantarà la resistència?}
Les resistències que fem servir al taller solen ser de \SI{0,25}{W} (un quart de watt). Una resistència de \SI{220}{\ohm} es connecta a \SI{5}{V}. L'aguantarà?
\begin{itemize}
 \item \textbf{Fórmula:} $P = \dfrac{V^2}{R}$
 \item \textbf{Resultat:} $P = \dfrac{5^2}{220} = \dfrac{25}{220} \approx \SI{0,114}{W}$
 \item \textbf{Conclusió:} $0{,}114 < 0{,}25$, per tant \textbf{sí}, l'aguanta.
\end{itemize}
\end{exemple}

\begin{atencio}
Calcular $I^2 \cdot R$ com si fos $(I \cdot R)^2$. Primer s'eleva al quadrat \textbf{només la intensitat} i després es multiplica per $R$.
\end{atencio}

\begin{guiada}{una resistència de \SI{50}{\ohm} amb \SI{0,2}{A}}
Quina potència dissipa? Aguantaria si fos de \SI{0,25}{W}?
\pas{1}{Dades:}{$R = \SI{50}{\ohm}$, \ $I = \SI{0,2}{A}$.}
\pas{2}{Fórmula (coneixem $I$ i $R$):}{\vspace{0.8cm}}
\pas{3}{Resultat:}{\vspace{0.8cm}}
\pas{4}{Aguanta \SI{0,25}{W}?}{\vspace{0.8cm}}
\resposta{0pt}{$P = I^2 R = 0{,}2^2 \cdot 50 = 0{,}04 \cdot 50 = \SI{2}{W}$. Com que $2 > 0{,}25$, \textbf{no} l'aguantaria: es cremaria.}
\end{guiada}

\subsection*{Exercicis}

\begin{exercici}{1}
Calcula la potència d'un aparell connectat a \SI{12}{V} i travessat per \SI{2}{A}.
\resposta{2.2cm}{$P = V \cdot I = 12 \cdot 2 = \SI{24}{W}$.}
\end{exercici}

\begin{exercici}{1}
Una pila de \SI{9}{V} fa circular \SI{0,1}{A} per un receptor. Quina potència té?
\resposta{2.2cm}{$P = 9 \cdot 0{,}1 = \SI{0,9}{W}$.}
\end{exercici}

\begin{exercici}{1}
Per una resistència de \SI{10}{\ohm} hi passa un corrent de \SI{0,5}{A}. Quina potència dissipa?
\resposta{2.2cm}{$P = I^2 R = 0{,}5^2 \cdot 10 = 0{,}25 \cdot 10 = \SI{2,5}{W}$.}
\end{exercici}

\begin{exercici}{2}
Una resistència d'\SI{1}{\kilo\ohm} es connecta a \SI{10}{V}. Quina potència dissipa? Aguantaria si fos de \SI{0,25}{W}?
\resposta{2.6cm}{$R = \SI{1000}{\ohm}$. $P = V^2/R = 100/1000 = \SI{0,1}{W}$. Com que $0{,}1 < 0{,}25$, \textbf{sí} l'aguantaria.}
\end{exercici}

\begin{exercici}{2}
Un aparell de \SI{60}{W} es connecta a \SI{12}{V}. Quina intensitat consumeix?\quad (Pista: aïlla $I$ de $P = V \cdot I$.)
\resposta{2.6cm}{$I = P/V = 60/12 = \SI{5}{A}$.}
\end{exercici}

\begin{exercici}{3}
Una resistència de \SI{100}{\ohm} i \SI{0,25}{W}. Quina és la \textbf{tensió màxima} que li podem aplicar sense passar-nos de la seva potència?\quad (Pista: aïlla $V$ de $P = V^2/R$.)
\resposta{3cm}{$V^2 = P \cdot R = 0{,}25 \cdot 100 = 25$, per tant $V = \sqrt{25} = \SI{5}{V}$.}
\end{exercici}

\begin{exercici}{3}
Amb la mateixa resistència de \SI{100}{\ohm} i \SI{0,25}{W}, quina és la \textbf{intensitat màxima} que hi pot passar?\quad Comprova que coincideix amb la tensió màxima de l'exercici anterior.
\resposta{3cm}{$I^2 = P/R = 0{,}25/100 = 0{,}0025$, per tant $I = \sqrt{0{,}0025} = \SI{0,05}{A} = \SI{50}{mA}$. Comprovació: $V = I \cdot R = 0{,}05 \cdot 100 = \SI{5}{V}$. \checkmark}
\end{exercici}

\begin{comprova}
Si dupliquem la tensió aplicada a una resistència, la potència que dissipa es fa el doble? Pensa en $P = V^2/R$.
\resposta{2cm}{No: es fa \textbf{quatre vegades} més gran, perquè la tensió està elevada al quadrat ($2^2 = 4$).}
\end{comprova}

\clearpage
\section{Resistències en sèrie}

Dues resistències estan \textbf{en sèrie} quan estan una darrere l'altra, en el mateix camí: el corrent no té on escollir i les travessa totes. Pensa en un tub d'aigua amb dues estretors seguides: tota l'aigua passa per les dues.

\begin{center}\cserie{$V_T$}{$R_1$}{$R_2$}\end{center}

\begin{idea}
En un circuit en sèrie:
\begin{enumerate}[label=\textbf{\arabic*.}]
 \item La \textbf{intensitat és la mateixa} a tot arreu: \quad $I_T = I_1 = I_2 = \cdots$
 \item Les \textbf{tensions se sumen}: \quad $V_T = V_1 + V_2 + \cdots$
 \item Les \textbf{resistències se sumen}: \quad $R_T = R_1 + R_2 + \cdots$
\end{enumerate}
\end{idea}

\begin{metode}{calcular un circuit en sèrie}
\begin{itemize}
 \item[\textbf{\ding{172}}] Suma les resistències: $R_T = R_1 + R_2 + \cdots$
 \item[\textbf{\ding{173}}] Calcula la intensitat total: $I_T = V_T / R_T$. Aquesta és la intensitat de \emph{cada} resistència.
 \item[\textbf{\ding{174}}] Calcula la tensió de cada resistència: $V_i = I \cdot R_i$.
 \item[\textbf{\ding{175}}] Calcula la potència de cada resistència: $P_i = V_i \cdot I$.
 \item[\textbf{\ding{176}}] \textbf{Comprova}: les tensions sumen $V_T$? les potències sumen $V_T \cdot I_T$?
\end{itemize}
\end{metode}

\begin{exemple}{dues resistències en sèrie}
Una pila de \SI{12}{V} alimenta dues resistències en sèrie: $R_1 = \SI{100}{\ohm}$ i $R_2 = \SI{200}{\ohm}$.
\begin{center}\cserie{\SI{12}{V}}{$R_1=\SI{100}{\ohm}$}{$R_2=\SI{200}{\ohm}$}\end{center}
\begin{itemize}
 \item $R_T = 100 + 200 = \SI{300}{\ohm}$
 \item $I_T = V_T / R_T = 12 / 300 = \SI{0,04}{A}$ \ (la mateixa per a $R_1$ i $R_2$)
 \item $V_1 = I \cdot R_1 = 0{,}04 \cdot 100 = \SI{4}{V}$ \qquad $V_2 = 0{,}04 \cdot 200 = \SI{8}{V}$
 \item $P_1 = 4 \cdot 0{,}04 = \SI{0,16}{W}$ \qquad $P_2 = 8 \cdot 0{,}04 = \SI{0,32}{W}$ \qquad $P_T = 12 \cdot 0{,}04 = \SI{0,48}{W}$
 \item \textbf{Comprovació:} $4 + 8 = 12$ \checkmark \qquad $0{,}16 + 0{,}32 = 0{,}48$ \checkmark
\end{itemize}
\begin{taulamag}
$R_1$ & 4 & 0{,}04 & 100 & 0{,}16\\
$R_2$ & 8 & 0{,}04 & 200 & 0{,}32\\
\midrule
Total & 12 & 0{,}04 & 300 & 0{,}48\\
\end{taulamag}
\end{exemple}

\begin{idea}
\textbf{Un truc útil: el divisor de tensió.} En sèrie, la tensió es reparteix \emph{en proporció a la resistència}. La resistència més gran se'n queda la part més gran:
\[ V_i = V_T \cdot \frac{R_i}{R_T} \]
A l'exemple: $V_1 = 12 \cdot \dfrac{100}{300} = \SI{4}{V}$.
\end{idea}

\begin{atencio}
Sumar les intensitats en un circuit en sèrie. \textbf{No!} La intensitat és \emph{la mateixa} a totes les resistències; el que se suma són les \textbf{tensions} i les \textbf{resistències}.
\end{atencio}

\begin{guiada}{pila de \SI{9}{V} amb $R_1 = \SI{300}{\ohm}$ i $R_2 = \SI{600}{\ohm}$ en sèrie}
\pas{1}{Resistència total:}{$R_T = $ \vspace{0.7cm}}
\pas{2}{Intensitat total:}{\vspace{0.7cm}}
\pas{3}{Tensions de cada resistència:}{\vspace{0.7cm}}
\begin{taulamag}
$R_1$ & & & 300 &\\
$R_2$ & & & 600 &\\
\midrule
Total & 9 & & &\\
\end{taulamag}
\resposta{0pt}{$R_T = 900\,\Omega$; $I = 9/900 = \SI{0,01}{A}$; $V_1 = 0{,}01 \cdot 300 = \SI{3}{V}$; $V_2 = 0{,}01 \cdot 600 = \SI{6}{V}$; $P_1 = \SI{0,03}{W}$, $P_2 = \SI{0,06}{W}$, $P_T = \SI{0,09}{W}$.
\begin{taulamag}
$R_1$ & 3 & 0{,}01 & 300 & 0{,}03\\
$R_2$ & 6 & 0{,}01 & 600 & 0{,}06\\
\midrule
Total & 9 & 0{,}01 & 900 & 0{,}09\\
\end{taulamag}}
\end{guiada}

\subsection*{Exercicis}

\begin{exercici}{1}
Calcula la resistència total de:\quad
\textbf{a)} \SI{100}{\ohm} i \SI{220}{\ohm} en sèrie.\qquad
\textbf{b)} \SI{1}{\kilo\ohm}, \SI{2,2}{\kilo\ohm} i \SI{4,7}{\kilo\ohm} en sèrie.
\resposta{2cm}{a) $100 + 220 = \SI{320}{\ohm}$. \quad b) $1 + 2{,}2 + 4{,}7 = \SI{7,9}{\kilo\ohm}$.}
\end{exercici}

\begin{exercici}[20]{1}
Completa la taula del circuit següent.
\begin{center}\cserie{\SI{9}{V}}{$R_1=\SI{100}{\ohm}$}{$R_2=\SI{200}{\ohm}$}\end{center}
\begin{taulamag}
$R_1$ & & & 100 &\\
$R_2$ & & & 200 &\\
\midrule
Total & 9 & & &\\
\end{taulamag}
\resposta{0pt}{$R_T = \SI{300}{\ohm}$, $I = \SI{0,03}{A}$.
\begin{taulamag}
$R_1$ & 3 & 0{,}03 & 100 & 0{,}09\\
$R_2$ & 6 & 0{,}03 & 200 & 0{,}18\\
\midrule
Total & 9 & 0{,}03 & 300 & 0{,}27\\
\end{taulamag}}
\end{exercici}

\begin{exercici}[20]{2}
Completa la taula del circuit següent.
\begin{center}\cserie{\SI{24}{V}}{$R_1=\SI{300}{\ohm}$}{$R_2=\SI{500}{\ohm}$}\end{center}
\begin{taulamag}
$R_1$ & & & 300 &\\
$R_2$ & & & 500 &\\
\midrule
Total & 24 & & &\\
\end{taulamag}
\resposta{0pt}{$R_T = \SI{800}{\ohm}$, $I = \SI{0,03}{A}$.
\begin{taulamag}
$R_1$ & 9 & 0{,}03 & 300 & 0{,}27\\
$R_2$ & 15 & 0{,}03 & 500 & 0{,}45\\
\midrule
Total & 24 & 0{,}03 & 800 & 0{,}72\\
\end{taulamag}}
\end{exercici}

\begin{exercici}[20]{2}
Completa la taula del circuit següent, amb tres resistències en sèrie.
\begin{center}\cserietres{\SI{12}{V}}{$R_1=\SI{100}{\ohm}$}{$R_2=\SI{200}{\ohm}$}{$R_3=\SI{300}{\ohm}$}\end{center}
\begin{taulamag}
$R_1$ & & & 100 &\\
$R_2$ & & & 200 &\\
$R_3$ & & & 300 &\\
\midrule
Total & 12 & & &\\
\end{taulamag}
\resposta{0pt}{$R_T = \SI{600}{\ohm}$, $I = \SI{0,02}{A}$.
\begin{taulamag}
$R_1$ & 2 & 0{,}02 & 100 & 0{,}04\\
$R_2$ & 4 & 0{,}02 & 200 & 0{,}08\\
$R_3$ & 6 & 0{,}02 & 300 & 0{,}12\\
\midrule
Total & 12 & 0{,}02 & 600 & 0{,}24\\
\end{taulamag}}
\end{exercici}

\begin{exercici}{2}
Dues resistències en sèrie, $R_1 = \SI{150}{\ohm}$ i una $R_2$ desconeguda, estan connectades a \SI{10}{V}. Per elles hi passa un corrent de \SI{0,05}{A}. Calcula $R_2$ i la tensió de cada resistència.
\resposta{3cm}{$R_T = V/I = 10/0{,}05 = \SI{200}{\ohm}$, per tant $R_2 = 200 - 150 = \SI{50}{\ohm}$. $V_1 = 0{,}05 \cdot 150 = \SI{7,5}{V}$, $V_2 = 0{,}05 \cdot 50 = \SI{2,5}{V}$. Comprovació: $7{,}5 + 2{,}5 = 10$ \checkmark}
\end{exercici}

\begin{exercici}{3}
Volem fer un \textbf{divisor de tensió}: una font de \SI{5}{V} alimenta $R_1 = \SI{300}{\ohm}$ i una $R_2$ en sèrie, i volem que entre els extrems de $R_2$ hi hagi exactament \SI{2}{V}. Quant ha de valer $R_2$?
\resposta{3.5cm}{$V_1 = 5 - 2 = \SI{3}{V}$. $I = V_1 / R_1 = 3/300 = \SI{0,01}{A}$. $R_2 = V_2 / I = 2 / 0{,}01 = \SI{200}{\ohm}$. Comprovació amb el divisor: $V_2 = 5 \cdot 200/500 = \SI{2}{V}$ \checkmark}
\end{exercici}

\begin{exercici}{3}
Una cadena de llums de Nadal té \num{20} bombetes iguals en sèrie i es connecta a \SI{230}{V}.
\textbf{a)} Quina tensió suporta cada bombeta?\quad \textbf{b)} Què passa amb la resta si una bombeta es fon? Per què?
\resposta{3cm}{a) $230 / 20 = \SI{11,5}{V}$ cada bombeta. b) S'apaguen totes: en fondre's una, el circuit s'obre i el corrent no pot circular enlloc.}
\end{exercici}
\clearpage
\section{Resistències en paral·lel}

Dues resistències estan \textbf{en paral·lel} quan els dos extrems de cadascuna estan connectats als mateixos dos punts: el corrent té \emph{dos camins} per escollir. Pensa en un tub d'aigua que es divideix en dos tubs i després torna a ajuntar-se: a més camins, més aigua pot passar.

\begin{center}\cparalel{$V_T$}{$R_1$}{$R_2$}\end{center}

\begin{idea}
En un circuit en paral·lel:
\begin{enumerate}[label=\textbf{\arabic*.}]
 \item La \textbf{tensió és la mateixa} a totes les branques: \quad $V_T = V_1 = V_2 = \cdots$
 \item La \textbf{intensitat es reparteix}: \quad $I_T = I_1 + I_2 + \cdots$
 \item La \textbf{resistència total} es calcula així: \quad $\dfrac{1}{R_T} = \dfrac{1}{R_1} + \dfrac{1}{R_2} + \cdots$
\end{enumerate}
\end{idea}

\begin{idea}
\textbf{Dos casos que s'usen molt:}
\begin{itemize}
 \item \textbf{Dues resistències:} $R_T = \dfrac{R_1 \cdot R_2}{R_1 + R_2}$
 \item \textbf{$n$ resistències iguals de valor $R$:} $R_T = \dfrac{R}{n}$
\end{itemize}
\textbf{Regla per comprovar:} la resistència total en paral·lel \textbf{sempre és més petita} que la més petita de les resistències.
\end{idea}

\begin{idea}
\textbf{Càlcul ràpid amb calculadora científica (tecla $x^{-1}$ o $1/x$):}
\\[0.3em]
Pots calcular la resistència total directament aïllant la fórmula com:
\[
  R_T = \left( \frac{1}{R_1} + \frac{1}{R_2} + \cdots \right)^{-1}
\]
\textbf{Passos a seguir:}
\begin{enumerate}[label=\textbf{\arabic*.}]
 \item Suma les inverses de cada resistència: \quad $R_1\ \mathbf{x^{-1}} + R_2\ \mathbf{x^{-1}} + \dots =$
 \item Inverteix el resultat obtingut: \quad $\mathbf{x^{-1}} =$ \quad (o bé $\text{Ans}^{-1} =$)
\end{enumerate}
\end{idea}

\begin{metode}{calcular un circuit en paral·lel}
\begin{itemize}
 \item\textbf{\ding{172}} Calcula $R_T$ amb la fórmula. \quad
 \item\textbf{\ding{173}} Calcula la intensitat total: $I_T = V_T / R_T$. \quad
 \item\textbf{\ding{174}} Calcula la intensitat de cada branca (la tensió és la mateixa): $I_i = V_T / R_i$. \quad
 \item\textbf{\ding{175}} Calcula les potències: $P_i = V_T \cdot I_i$. \quad
 \item\textbf{\ding{176}} \textbf{Comprova}: les intensitats sumen $I_T$? les potències sumen $V_T \cdot I_T$?
\end{itemize}
\end{metode}

\begin{exemple}{dues resistències en paral·lel}
Una pila de \SI{12}{V} alimenta $R_1 = \SI{200}{\ohm}$ i $R_2 = \SI{300}{\ohm}$ en paral·lel.
\begin{center}\cparalel{\SI{12}{V}}{$R_1=\SI{200}{\ohm}$}{$R_2=\SI{300}{\ohm}$}\end{center}
\begin{itemize}
 \item $R_T = \dfrac{200 \cdot 300}{200 + 300} = \dfrac{60000}{500} = \SI{120}{\ohm}$ \ (menor que \SI{200}{\ohm}: té sentit)
 \item $I_T = 12 / 120 = \SI{0,1}{A}$
 \item $I_1 = 12 / 200 = \SI{0,06}{A}$ \qquad $I_2 = 12 / 300 = \SI{0,04}{A}$
 \item $P_1 = 12 \cdot 0{,}06 = \SI{0,72}{W}$ \qquad $P_2 = 12 \cdot 0{,}04 = \SI{0,48}{W}$ \qquad $P_T = 12 \cdot 0{,}1 = \SI{1,2}{W}$
 \item \textbf{Comprovació:} $0{,}06 + 0{,}04 = 0{,}1$ \checkmark \qquad $0{,}72 + 0{,}48 = 1{,}2$ \checkmark
\end{itemize}
\begin{taulamag}
$R_1$ & 12 & 0{,}06 & 200 & 0{,}72\\
$R_2$ & 12 & 0{,}04 & 300 & 0{,}48\\
\midrule
Total & 12 & 0{,}1 & 120 & 1{,}2\\
\end{taulamag}
\end{exemple}

\begin{exemple}{resistències iguals}
Tres resistències de \SI{90}{\ohm} en paral·lel: $R_T = \dfrac{90}{3} = \SI{30}{\ohm}$.
\end{exemple}

\begin{atencio}
\begin{itemize}
 \item \textbf{(1)} Sumar les resistències directament, com en sèrie. En paral·lel \emph{no} es sumen. \quad
 \item \textbf{(2)} Oblidar l'últim pas: si has obtingut $\dfrac{1}{R_T} = 0{,}05$, encara no has acabat: $R_T = \dfrac{1}{0{,}05} = \SI{20}{\ohm}$.
\end{itemize}
\end{atencio}

\clearpage
\begin{guiada}{pila de \SI{6}{V} amb $R_1 = \SI{100}{\ohm}$ i $R_2 = \SI{300}{\ohm}$ en paral·lel}
\pas{1}{Resistència total:}{\vspace{0.7cm}}
\pas{2}{Intensitat total:}{\vspace{0.7cm}}
\pas{3}{Intensitat de cada branca:}{\vspace{0.7cm}}
\begin{taulamag}
$R_1$ & 6 & & 100 &\\
$R_2$ & 6 & & 300 &\\
\midrule
Total & 6 & & &\\
\end{taulamag}
\resposta{0pt}{$R_T = \dfrac{100 \cdot 300}{400} = \SI{75}{\ohm}$; $I_T = 6/75 = \SI{0,08}{A}$; $I_1 = 6/100 = \SI{0,06}{A}$; $I_2 = 6/300 = \SI{0,02}{A}$.
\begin{taulamag}
$R_1$ & 6 & 0{,}06 & 100 & 0{,}36\\
$R_2$ & 6 & 0{,}02 & 300 & 0{,}12\\
\midrule
Total & 6 & 0{,}08 & 75 & 0{,}48\\
\end{taulamag}}
\end{guiada}
\clearpage
\subsection*{Exercicis}

\begin{exercici}{1}
Calcula la resistència total en paral·lel de:\quad
\textbf{a)} dues resistències de \SI{100}{\ohm}.\qquad
\textbf{b)} tres resistències de \SI{60}{\ohm}.\qquad
\textbf{c)} \SI{40}{\ohm} i \SI{60}{\ohm}.
\resposta{2.4cm}{a) $100/2 = \SI{50}{\ohm}$. \quad b) $60/3 = \SI{20}{\ohm}$. \quad c) $\dfrac{40 \cdot 60}{100} = \SI{24}{\ohm}$.}
\end{exercici}

\begin{exercici}{1}
Calcula la resistència total en paral·lel de:\quad
\textbf{a)} \SI{12}{\ohm} i \SI{4}{\ohm}.\qquad
\textbf{b)} \SI{10}{\ohm} i \SI{40}{\ohm}.
\resposta{2.4cm}{a) $\dfrac{12 \cdot 4}{16} = \SI{3}{\ohm}$. \quad b) $\dfrac{10 \cdot 40}{50} = \SI{8}{\ohm}$.}
\end{exercici}

\begin{exercici}[20]{2}
Completa la taula del circuit següent.
\begin{center}\cparalel{\SI{12}{V}}{$R_1=\SI{40}{\ohm}$}{$R_2=\SI{60}{\ohm}$}\end{center}
\begin{taulamag}
$R_1$ & 12 & & 40 &\\
$R_2$ & 12 & & 60 &\\
\midrule
Total & 12 & & &\\
\end{taulamag}
\resposta{0pt}{$R_T = \SI{24}{\ohm}$, $I_T = \SI{0,5}{A}$.
\begin{taulamag}
$R_1$ & 12 & 0{,}3 & 40 & 3{,}6\\
$R_2$ & 12 & 0{,}2 & 60 & 2{,}4\\
\midrule
Total & 12 & 0{,}5 & 24 & 6\\
\end{taulamag}}
\end{exercici}

\begin{exercici}[20]{2}
Completa la taula del circuit següent, amb tres resistències en paral·lel.\quad (Pista: $\frac{1}{R_T} = \frac{1}{100} + \frac{1}{200} + \frac{1}{200}$.)
\begin{center}\cparaleltres{\SI{10}{V}}{$R_1=\SI{100}{\ohm}$}{$R_2=\SI{200}{\ohm}$}{$R_3=\SI{200}{\ohm}$}\end{center}
\begin{taulamag}
$R_1$ & 10 & & 100 &\\
$R_2$ & 10 & & 200 &\\
$R_3$ & 10 & & 200 &\\
\midrule
Total & 10 & & &\\
\end{taulamag}
\resposta{0pt}{$\dfrac{1}{R_T} = 0{,}01 + 0{,}005 + 0{,}005 = 0{,}02$, per tant $R_T = \SI{50}{\ohm}$; $I_T = \SI{0,2}{A}$.
\begin{taulamag}
$R_1$ & 10 & 0{,}1 & 100 & 1\\
$R_2$ & 10 & 0{,}05 & 200 & 0{,}5\\
$R_3$ & 10 & 0{,}05 & 200 & 0{,}5\\
\midrule
Total & 10 & 0{,}2 & 50 & 2\\
\end{taulamag}}
\end{exercici}

\begin{exercici}{3}
Una resistència de \SI{60}{\ohm} es connecta en paral·lel amb una altra resistència $R_2$, i la resistència total és de \SI{40}{\ohm}. Quant val $R_2$?
\resposta{3cm}{$\dfrac{1}{R_2} = \dfrac{1}{40} - \dfrac{1}{60} = \dfrac{3 - 2}{120} = \dfrac{1}{120}$, per tant $R_2 = \SI{120}{\ohm}$. Comprovació: $\dfrac{60 \cdot 120}{180} = \SI{40}{\ohm}$ \checkmark}
\end{exercici}

\begin{exercici}{3}
Dues resistències $R_1$ i $R_2$ en paral·lel estan connectades a \SI{6}{V}. Un amperímetre indica que la intensitat total és \SI{0,5}{A} i que per $R_1$ hi passa \SI{0,2}{A}. Calcula $I_2$, $R_1$, $R_2$ i $R_T$.
\resposta{3.5cm}{$I_2 = 0{,}5 - 0{,}2 = \SI{0,3}{A}$. $R_1 = 6/0{,}2 = \SI{30}{\ohm}$. $R_2 = 6/0{,}3 = \SI{20}{\ohm}$. $R_T = 6/0{,}5 = \SI{12}{\ohm}$. Comprovació: $\dfrac{30 \cdot 20}{50} = \SI{12}{\ohm}$ \checkmark}
\end{exercici}

\begin{comprova}
A casa, els endolls estan connectats en paral·lel. Per què? Què passaria si estiguessin en sèrie?
\resposta{2.4cm}{En paral·lel, cada aparell rep la tensió completa (\SI{230}{V}) i funciona independentment dels altres. En sèrie, la tensió es repartiria i, si un aparell s'apagués, s'apagarien tots.}
\end{comprova}

\clearpage
\section{Circuits mixtos}

Un circuit \textbf{mixt} combina resistències en sèrie i en paral·lel. Sembla més difícil, però no ho és: es resol amb un \textbf{mètode fix} que usa només el que ja saps.

\begin{center}
\cmixa{$V_T$}{$R_1$}{$R_2$}{$R_3$}\qquad\cmixc{$V_T$}{$R_1$}{$R_2$}{$R_3$}
\end{center}
\begin{center}\small A l'esquerra, $R_1$ en sèrie amb ($R_2$ en paral·lel amb $R_3$). A la dreta, ($R_1$ en sèrie amb $R_2$) en paral·lel amb $R_3$.\end{center}

\begin{idea}
La clau és \textbf{reduir} i després \textbf{desfer}:
\begin{itemize}
 \item \textbf{Reduir:} substitueix cada grup per una sola resistència equivalent, fins que el circuit sigui una sola resistència total.
 \item \textbf{Desfer:} tornes enrere pas a pas, calculant les tensions i les intensitats de cada grup i, dins de cada grup, de cada resistència.
\end{itemize}
\end{idea}

\begin{metode}{circuit mixt (reduir i desfer)}
\begin{itemize}
 \item \textbf{\ding{172}} Busca els \textbf{grups} en paral·lel (o en sèrie) i calcula la resistència equivalent de cada un. \quad
 \item \textbf{\ding{173}} \textbf{Redibuixa} el circuit amb els equivalents i calcula $R_T$. \quad
 \item \textbf{\ding{174}} Calcula la intensitat total: $I_T = V_T / R_T$. \quad
 \item \textbf{\ding{175}} \textbf{Desfés:} en les parts en sèrie, la intensitat és la mateixa; calcula la tensió de cada grup amb $V = I \cdot R_{equivalent}$. \quad
 \item \textbf{\ding{176}} Dins de cada grup en paral·lel, la \textbf{tensió és la del grup}: calcula la intensitat de cada branca amb $I = V / R$. \quad
 \item \textbf{\ding{177}} Omple la taula i \textbf{comprova}: les potències sumen $V_T \cdot I_T$?
\end{itemize}
\end{metode}

\begin{exemple}{$R_1$ en sèrie amb ($R_2$ en paral·lel amb $R_3$)}
$V_T = \SI{12}{V}$, $R_1 = \SI{30}{\ohm}$, $R_2 = \SI{60}{\ohm}$, $R_3 = \SI{60}{\ohm}$.
\begin{center}\cmixa{\SI{12}{V}}{$R_1=\SI{30}{\ohm}$}{$R_2=\SI{60}{\ohm}$}{$R_3=\SI{60}{\ohm}$}\end{center}
\begin{itemize}
 \item \textbf{\ding{172} Reduir el grup en paral·lel:} $R_{23} = \dfrac{60 \cdot 60}{120} = \SI{30}{\ohm}$.
 \item \textbf{\ding{173} Circuit simplificat} (ara tot és en sèrie): $R_T = R_1 + R_{23} = 30 + 30 = \SI{60}{\ohm}$.
\begin{center}\cserie{\SI{12}{V}}{$R_1=\SI{30}{\ohm}$}{$R_{23}=\SI{30}{\ohm}$}\end{center}
 \item \textbf{\ding{174}} $I_T = 12 / 60 = \SI{0,2}{A}$.
 \item \textbf{\ding{175} Desfer:} $V_1 = 0{,}2 \cdot 30 = \SI{6}{V}$ \ i \ $V_{23} = 0{,}2 \cdot 30 = \SI{6}{V}$ \ (comprovació: $6 + 6 = 12$ \checkmark).
 \item \textbf{\ding{176}} Dins del grup, la tensió és \SI{6}{V}: $I_2 = 6/60 = \SI{0,1}{A}$ \ i \ $I_3 = 6/60 = \SI{0,1}{A}$ (suma \SI{0,2}{A} \checkmark).
\end{itemize}
\begin{taulamag}
$R_1$ & 6 & 0{,}2 & 30 & 1{,}2\\
$R_2$ & 6 & 0{,}1 & 60 & 0{,}6\\
$R_3$ & 6 & 0{,}1 & 60 & 0{,}6\\
\midrule
Total & 12 & 0{,}2 & 60 & 2{,}4\\
\end{taulamag}
\textbf{Comprovació de potències:} $1{,}2 + 0{,}6 + 0{,}6 = 2{,}4 = 12 \cdot 0{,}2$ \checkmark
\end{exemple}

\begin{exemple}{($R_1$ en sèrie amb $R_2$) en paral·lel amb $R_3$}
$V_T = \SI{12}{V}$, $R_1 = \SI{100}{\ohm}$, $R_2 = \SI{200}{\ohm}$, $R_3 = \SI{150}{\ohm}$.
\begin{center}\cmixc{\SI{12}{V}}{$R_1=\SI{100}{\ohm}$}{$R_2=\SI{200}{\ohm}$}{$R_3=\SI{150}{\ohm}$}\end{center}
\begin{itemize}
 \item \textbf{\ding{172} Reduir el grup en sèrie:} $R_{12} = 100 + 200 = \SI{300}{\ohm}$.
 \item \textbf{\ding{173} Circuit simplificat} (ara són dues en paral·lel): $R_T = \dfrac{300 \cdot 150}{300 + 150} = \dfrac{45000}{450} = \SI{100}{\ohm}$.
 \item \textbf{\ding{174}} $I_T = 12 / 100 = \SI{0,12}{A}$.
 \item \textbf{\ding{175} Desfer:} les dues branques tenen \SI{12}{V}. Branca de $R_{12}$: $I_{12} = 12/300 = \SI{0,04}{A}$. Branca de $R_3$: $I_3 = 12/150 = \SI{0,08}{A}$ (suma \SI{0,12}{A} \checkmark).
 \item \textbf{\ding{176}} A la branca en sèrie, la intensitat és la mateixa: $V_1 = 0{,}04 \cdot 100 = \SI{4}{V}$ \ i \ $V_2 = 0{,}04 \cdot 200 = \SI{8}{V}$ (suma \SI{12}{V} \checkmark).
\end{itemize}
\begin{taulamag}
$R_1$ & 4 & 0{,}04 & 100 & 0{,}16\\
$R_2$ & 8 & 0{,}04 & 200 & 0{,}32\\
$R_3$ & 12 & 0{,}08 & 150 & 0{,}96\\
\midrule
Total & 12 & 0{,}12 & 100 & 1{,}44\\
\end{taulamag}
\textbf{Comprovació de potències:} $0{,}16 + 0{,}32 + 0{,}96 = 1{,}44 = 12 \cdot 0{,}12$ \checkmark
\end{exemple}

\begin{atencio}
\textbf{(1)} Sumar totes les resistències com si fossin en sèrie. Primer cal reduir els grups en paral·lel. \quad
\textbf{(2)} Oblidar que dins d'un grup en paral·lel \textbf{la tensió és la mateixa} per a totes les branques i que la \textbf{intensitat es reparteix}.
\end{atencio}

\begin{guiada}{$R_1$ en sèrie amb ($R_2$ en paral·lel amb $R_3$)}
$V_T = \SI{10}{V}$, $R_1 = \SI{50}{\ohm}$, $R_2 = \SI{100}{\ohm}$, $R_3 = \SI{100}{\ohm}$.
\begin{center}\cmixa{\SI{10}{V}}{$R_1=\SI{50}{\ohm}$}{$R_2=\SI{100}{\ohm}$}{$R_3=\SI{100}{\ohm}$}\end{center}
\pas{1}{Resistència equivalent del grup en paral·lel $R_{23}$:}{\vspace{0.7cm}}
\pas{2}{Resistència total $R_T$ i intensitat total $I_T$:}{\vspace{0.7cm}}
\pas{3}{Tensió de $R_1$ i del grup $R_{23}$:}{\vspace{0.7cm}}
\pas{4}{Intensitat de $R_2$ i de $R_3$:}{\vspace{0.7cm}}
\begin{taulamag}
$R_1$ & & & 50 &\\
$R_2$ & & & 100 &\\
$R_3$ & & & 100 &\\
\midrule
Total & 10 & & &\\
\end{taulamag}
\resposta{0pt}{$R_{23} = \SI{50}{\ohm}$; $R_T = 100\,\Omega$ i $I_T = \SI{0,1}{A}$; $V_1 = \SI{5}{V}$ i $V_{23} = \SI{5}{V}$; $I_2 = I_3 = \SI{0,05}{A}$.
\begin{taulamag}
$R_1$ & 5 & 0{,}1 & 50 & 0{,}5\\
$R_2$ & 5 & 0{,}05 & 100 & 0{,}25\\
$R_3$ & 5 & 0{,}05 & 100 & 0{,}25\\
\midrule
Total & 10 & 0{,}1 & 100 & 1\\
\end{taulamag}}
\end{guiada}

\subsection*{Exercicis}

\begin{exercici}{1}
Calcula la resistència total d'aquests dos circuits.
\begin{center}
\cmixa{\SI{12}{V}}{$R_1=\SI{10}{\ohm}$}{$R_2=\SI{40}{\ohm}$}{$R_3=\SI{40}{\ohm}$}\qquad
\cmixc{\SI{12}{V}}{$R_1=\SI{50}{\ohm}$}{$R_2=\SI{50}{\ohm}$}{$R_3=\SI{100}{\ohm}$}
\end{center}
\resposta{0pt}{Esquerra: $R_{23} = 40/2 = \SI{20}{\ohm}$, $R_T = 10 + 20 = \SI{30}{\ohm}$. Dreta: $R_{12} = 50 + 50 = \SI{100}{\ohm}$, $R_T = \dfrac{100 \cdot 100}{200} = \SI{50}{\ohm}$.}
\end{exercici}

\begin{exercici}[20]{2}
Completa la taula del circuit següent.
\begin{center}\cmixa{\SI{24}{V}}{$R_1=\SI{60}{\ohm}$}{$R_2=\SI{120}{\ohm}$}{$R_3=\SI{120}{\ohm}$}\end{center}
\begin{taulamag}
$R_1$ & & & 60 &\\
$R_2$ & & & 120 &\\
$R_3$ & & & 120 &\\
\midrule
Total & 24 & & &\\
\end{taulamag}
\resposta{0pt}{$R_{23} = \SI{60}{\ohm}$, $R_T = \SI{120}{\ohm}$, $I_T = \SI{0,2}{A}$.
\begin{taulamag}
$R_1$ & 12 & 0{,}2 & 60 & 2{,}4\\
$R_2$ & 12 & 0{,}1 & 120 & 1{,}2\\
$R_3$ & 12 & 0{,}1 & 120 & 1{,}2\\
\midrule
Total & 24 & 0{,}2 & 120 & 4{,}8\\
\end{taulamag}}
\end{exercici}

\begin{exercici}[20]{2}
Completa la taula del circuit següent.
\begin{center}\cmixc{\SI{10}{V}}{$R_1=\SI{20}{\ohm}$}{$R_2=\SI{30}{\ohm}$}{$R_3=\SI{50}{\ohm}$}\end{center}
\begin{taulamag}
$R_1$ & & & 20 &\\
$R_2$ & & & 30 &\\
$R_3$ & & & 50 &\\
\midrule
Total & 10 & & &\\
\end{taulamag}
\resposta{0pt}{$R_{12} = \SI{50}{\ohm}$, $R_T = \dfrac{50 \cdot 50}{100} = \SI{25}{\ohm}$, $I_T = \SI{0,4}{A}$.
\begin{taulamag}
$R_1$ & 4 & 0{,}2 & 20 & 0{,}8\\
$R_2$ & 6 & 0{,}2 & 30 & 1{,}2\\
$R_3$ & 10 & 0{,}2 & 50 & 2\\
\midrule
Total & 10 & 0{,}4 & 25 & 4\\
\end{taulamag}}
\end{exercici}

\begin{exercici}[20]{3}
Completa la taula del circuit següent, amb quatre resistències.\quad (Pista: comença pel grup en paral·lel.)
\begin{center}\cmixd{\SI{12}{V}}{$R_1=\SI{20}{\ohm}$}{$R_2=\SI{60}{\ohm}$}{$R_3=\SI{60}{\ohm}$}{$R_4=\SI{10}{\ohm}$}\end{center}
\begin{taulamag}
$R_1$ & & & 20 &\\
$R_2$ & & & 60 &\\
$R_3$ & & & 60 &\\
$R_4$ & & & 10 &\\
\midrule
Total & 12 & & &\\
\end{taulamag}
\resposta{0pt}{$R_{23} = \SI{30}{\ohm}$, $R_T = 20 + 30 + 10 = \SI{60}{\ohm}$, $I_T = \SI{0,2}{A}$.
\begin{taulamag}
$R_1$ & 4 & 0{,}2 & 20 & 0{,}8\\
$R_2$ & 6 & 0{,}1 & 60 & 0{,}6\\
$R_3$ & 6 & 0{,}1 & 60 & 0{,}6\\
$R_4$ & 2 & 0{,}2 & 10 & 0{,}4\\
\midrule
Total & 12 & 0{,}2 & 60 & 2{,}4\\
\end{taulamag}}
\end{exercici}

\begin{exercici}{3}
En el circuit ($R_1$ en sèrie amb $R_2$ en paral·lel amb $R_3$) de la figura, $V_T = \SI{12}{V}$, $R_2 = R_3 = \SI{200}{\ohm}$ i l'amperímetre indica $I_T = \SI{0,1}{A}$. Calcula $R_1$ i la tensió de $R_1$.
\begin{center}\cmixa{\SI{12}{V}}{$R_1=?$}{$R_2=\SI{200}{\ohm}$}{$R_3=\SI{200}{\ohm}$}\end{center}
\resposta{3.5cm}{$R_T = 12/0{,}1 = \SI{120}{\ohm}$. $R_{23} = 200/2 = \SI{100}{\ohm}$. Per tant $R_1 = 120 - 100 = \SI{20}{\ohm}$ i $V_1 = 0{,}1 \cdot 20 = \SI{2}{V}$ (i $V_{23} = \SI{10}{V}$; $2 + 10 = 12$ \checkmark).}
\end{exercici}

\begin{exercici}{3}
\textbf{Troba l'error.} Per al circuit de l'exemple ($V_T = \SI{12}{V}$, $R_1 = \SI{30}{\ohm}$, $R_2 = R_3 = \SI{60}{\ohm}$ amb $R_2$ i $R_3$ en paral·lel), un company ha calculat $R_T = 30 + 60 + 60 = \SI{150}{\ohm}$. Quin error ha fet? Quina és la resistència total correcta?
\resposta{3cm}{Ha sumat les tres resistències com si fossin totes en sèrie. Però $R_2$ i $R_3$ estan en paral·lel: $R_{23} = 30\,\Omega$, i $R_T = 30 + 30 = \SI{60}{\ohm}$.}
\end{exercici}
\clearpage
\section{Repte final}

Ja saps calcular circuits en sèrie, en paral·lel i mixtos. Aquí tens tres reptes per posar-ho tot en pràctica. Alguns són de \textbf{dissenyar}, no només de calcular.

\begin{exercici}{3}
\textbf{El repte de les resistències de \SI{100}{\ohm}.} Necessites una resistència de \SI{150}{\ohm}, però només tens resistències de \SI{100}{\ohm} (en tens moltes). Com ho pots fer? Dibuixa el muntatge i comprova que la resistència total és \SI{150}{\ohm}.
\resposta{4cm}{Una resistència de \SI{100}{\ohm} \textbf{en sèrie} amb dues de \SI{100}{\ohm} \textbf{en paral·lel} (que fan \SI{50}{\ohm}): $100 + 50 = \SI{150}{\ohm}$.
\begin{center}\cmixa{\SI{9}{V}}{$R_1=\SI{100}{\ohm}$}{$R_2=\SI{100}{\ohm}$}{$R_3=\SI{100}{\ohm}$}\end{center}}
\end{exercici}

\begin{exercici}{3}
\textbf{El divisor de tensió.} Tens una pila de \SI{9}{V} i necessites obtenir \SI{3}{V} entre els extrems d'una resistència $R_2$, connectada en sèrie amb $R_1 = \SI{200}{\ohm}$. Quant ha de valer $R_2$? Comprova-ho calculant la intensitat i la tensió de cada resistència.
\resposta{3.5cm}{$V_1 = 9 - 3 = \SI{6}{V}$. $I = 6/200 = \SI{0,03}{A}$. $R_2 = 3/0{,}03 = \SI{100}{\ohm}$. Comprovació: $R_T = 300\,\Omega$, $I = 9/300 = \SI{0,03}{A}$, $V_2 = 0{,}03 \cdot 100 = \SI{3}{V}$ \checkmark}
\end{exercici}

\begin{exercici}{3}
\textbf{El detectiu de circuits.} Has connectat dues resistències $R_1$ i $R_2$ a una font de \SI{12}{V} i has fet unes mesures amb el polímetre. Per a cada cas, digues si estan \textbf{en sèrie o en paral·lel} i calcula $R_1$ i $R_2$.
\begin{center}
\renewcommand{\arraystretch}{1.5}
\begin{tabular}{@{}lcccc@{}}
\toprule
 & \textbf{Tensió de $R_1$} & \textbf{Tensió de $R_2$} & \textbf{Intensitat per $R_1$} & \textbf{Intensitat per $R_2$}\\
\midrule
Cas A & \SI{4}{V} & \SI{8}{V} & \SI{40}{mA} & \SI{40}{mA}\\
Cas B & \SI{12}{V} & \SI{12}{V} & \SI{60}{mA} & \SI{40}{mA}\\
\bottomrule
\end{tabular}
\end{center}
\resposta{4cm}{\textbf{Cas A: sèrie.} La intensitat és la mateixa i les tensions sumen $4 + 8 = 12$. $R_1 = 4/0{,}04 = \SI{100}{\ohm}$ i $R_2 = 8/0{,}04 = \SI{200}{\ohm}$.\\
\textbf{Cas B: paral·lel.} La tensió és la mateixa i les intensitats se sumen ($I_T = \SI{100}{mA}$). $R_1 = 12/0{,}06 = \SI{200}{\ohm}$ i $R_2 = 12/0{,}04 = \SI{300}{\ohm}$.}
\end{exercici}

\begin{exercici}{2}
\textbf{Resum.} Completa la taula amb el que has après sobre sèrie i paral·lel.
\begin{center}
\renewcommand{\arraystretch}{2.2}
\begin{tabular}{@{}l>{\centering\arraybackslash}p{5.5cm}>{\centering\arraybackslash}p{5.5cm}@{}}
\toprule
 & \textbf{Sèrie} & \textbf{Paral·lel}\\
\midrule
\textbf{Intensitat} & & \\
\textbf{Tensió} & & \\
\textbf{Resistència total} & & \\
\textbf{Si una es fon\ldots} & & \\
\bottomrule
\end{tabular}
\end{center}
\resposta{0pt}{
\begin{center}
\renewcommand{\arraystretch}{1.5}
\begin{tabular}{@{}l>{\centering\arraybackslash}p{5.5cm}>{\centering\arraybackslash}p{5.5cm}@{}}
\toprule
 & \textbf{Sèrie} & \textbf{Paral·lel}\\
\midrule
\textbf{Intensitat} & és la mateixa & es reparteix: $I_T = I_1 + I_2$\\
\textbf{Tensió} & es reparteix: $V_T = V_1 + V_2$ & és la mateixa\\
\textbf{Resistència total} & $R_T = R_1 + R_2$ & $\frac{1}{R_T} = \frac{1}{R_1} + \frac{1}{R_2}$ (menor que la menor)\\
\textbf{Si una es fon\ldots} & s'apaguen totes & les altres continuen\\
\bottomrule
\end{tabular}
\end{center}}
\end{exercici}

\subsection*{Autoavaluació}

Marca el que ja saps fer sense ajuda:

\begin{itemize}[label=$\square$, itemsep=0.4em]
 \item Passar mA, k$\Omega$, M$\Omega$\ldots a unitats base i tornar enrere.
 \item Aplicar la llei d'Ohm per trobar $V$, $I$ o $R$.
 \item Calcular la potència d'una resistència i saber si l'aguanta.
 \item Calcular un circuit en sèrie i omplir la taula de magnituds.
 \item Calcular un circuit en paral·lel i omplir la taula de magnituds.
 \item Reduir i desfer un circuit mixt.
 \item Comprovar els resultats (les tensions, les intensitats i les potències sumen bé).
\end{itemize}

\clearpage
\section*{Full de fórmules}
\addcontentsline{toc}{section}{Full de fórmules}

\begin{tcolorbox}[colback=blau!5, colframe=blau, title={Llei d'Ohm i potència}]
\[ V = I \cdot R \qquad I = \frac{V}{R} \qquad R = \frac{V}{I} \]
\[ P = V \cdot I \qquad P = I^2 \cdot R \qquad P = \frac{V^2}{R} \]
\end{tcolorbox}

\begin{tcolorbox}[colback=verd!6, colframe=verd, title={En sèrie}]
La \textbf{intensitat és la mateixa}; les \textbf{tensions se sumen}.
\[ R_T = R_1 + R_2 + \cdots \qquad V_T = V_1 + V_2 + \cdots \qquad V_i = V_T \cdot \frac{R_i}{R_T} \]
\end{tcolorbox}

\begin{tcolorbox}[colback=taronja!8, colframe=taronja, title={En paral·lel}]
La \textbf{tensió és la mateixa}; les \textbf{intensitats se sumen}.
\[ \frac{1}{R_T} = \frac{1}{R_1} + \frac{1}{R_2} + \cdots \qquad I_T = I_1 + I_2 + \cdots \]
\[ \text{Dues resistències: } R_T = \frac{R_1 \cdot R_2}{R_1 + R_2} \qquad\quad n \text{ iguals: } R_T = \frac{R}{n} \]
\end{tcolorbox}

\begin{tcolorbox}[colback=lila!6, colframe=lila, title={Circuit mixt: reduir i desfer}]
\textbf{\ding{172}} Calcula els equivalents dels grups. \ \textbf{\ding{173}} Redibuixa i calcula $R_T$. \ \textbf{\ding{174}} $I_T = V_T / R_T$. \ \textbf{\ding{175}} Desfés: tensió de cada grup. \ \textbf{\ding{176}} Intensitat de cada resistència. \ \textbf{\ding{177}} Comprova les potències.
\end{tcolorbox}

\begin{center}
\renewcommand{\arraystretch}{1.4}
\begin{tabular}{@{}llll@{}}
\toprule
\textbf{Prefix} & \textbf{Símbol} & \textbf{Factor} & \textbf{Per passar a unitat base}\\
\midrule
micro & $\mu$ & $10^{-6}$ & dividir per $1\,000\,000$\\
mil·li & m & $10^{-3}$ & dividir per $1000$\\
quilo & k & $10^{3}$ & multiplicar per $1000$\\
mega & M & $10^{6}$ & multiplicar per $1\,000\,000$\\
\bottomrule
\end{tabular}
\end{center}


\end{document}
